\section{Relational Databases}

% ─────────────────────────────────────────────────────────────────────────────
\begin{frame}{Edgar Frank ``Ted'' Codd}
  \begin{block}{Biographical Overview}
    \begin{itemize}
      \item \textbf{19 August 1923 -- 18 April 2003}
      \item Inventor of the \textbf{relational data model}, the concept of
            \textbf{OLAP} (Online Analytical Processing), and co-designer of
            \textbf{SEQUEL} (the precursor to SQL).
    \end{itemize}
  \end{block}

  \begin{block}{Career Timeline at IBM}
    \begin{itemize}
      \item \textbf{1949}: Joins IBM New York as a developer.
      \item \textbf{1965}: PhD in computer science, University of Michigan
            (funded by IBM scholarship).
      \item \textbf{1967}: Moves to IBM Almaden Research Center, San José.
      \item \textbf{1969--1970}: Publishes the foundational papers on the
            relational model, normal forms, and a relational sub-language.
      \item \textbf{1970s}: Co-develops SEQUEL with Donald Chamberlin.
    \end{itemize}
  \end{block}

  Codd's work matters because \textbf{almost all modern business software} --
  banking, e-commerce, ERP, CRM, logistics -- runs on relational databases
  that implement his mathematical framework.  He won the \textbf{Turing Award}
  in 1981 for his contributions, the highest honour in computer science.
\end{frame}

% ─────────────────────────────────────────────────────────────────────────────
\begin{frame}{Relations}
  \begin{block}{Definition}
    A \textbf{relation} is a mathematical set of tuples that all share the same
    structure (schema).  It can be visualised as a \textbf{table} with named
    columns and a collection of rows.
    \begin{itemize}
      \item Columns = \textbf{Attributes} (named, typed)
      \item Rows = \textbf{Tuples} (concrete data records)
      \item A table contains $0 \ldots *$ tuples.
    \end{itemize}
  \end{block}

  \begin{examples}{}
    \texttt{Album(\underline{ID}, EAN, Title, Main\_Artist, Release\_Date)}\\
    \texttt{Track(\underline{ID}, Album\_ID\textsuperscript{FK}, Disc\_Nr, Track\_Nr, ISRC, Title, Duration)}
  \end{examples}

  The word \emph{relation} comes from the Latin \emph{relatio} -- the
  relationship between entities of the same kind.  Because a relation is a
  \textbf{mathematical set}, it contains \emph{no duplicate tuples} in the
  pure theoretical definition.  In practice, SQL allows duplicates unless you
  explicitly use \texttt{DISTINCT} or enforce a uniqueness constraint -- an
  important distinction between the theory and its implementation.

  The \textbf{closure property} of the relational model states that the
  result of any relational operation is itself a relation, which means
  operations can be arbitrarily composed and nested.
\end{frame}

% ─────────────────────────────────────────────────────────────────────────────
\begin{frame}{Attributes}
  \begin{block}{Definition}
    An \textbf{attribute} is a named column in a relation that defines one
    specific piece of data about the entities stored in that relation.
    Every attribute has:
    \begin{itemize}
      \item A \textbf{name} (e.g.\ \texttt{Release\_Date})
      \item A \textbf{domain} -- the set of all allowed values
            (e.g.\ \texttt{DATE}, \texttt{VARCHAR(255)}, \texttt{INTEGER})
    \end{itemize}
  \end{block}

  \begin{examples}{}
    Attributes of \texttt{Album}:
    \texttt{ID}, \texttt{EAN}, \texttt{Title}, \texttt{Main\_Artist},
    \texttt{Release\_Date}\\[4pt]
    The domain of \texttt{EAN} is a 13-digit integer;
    the domain of \texttt{Title} is a non-empty string of up to 255 characters.
  \end{examples}

  Attributes describe the \emph{properties} or \emph{characteristics} of the
  entities stored in a table.  Choosing good attribute names is not a minor
  stylistic concern -- unclear names lead to misunderstandings, bugs, and
  costly data migrations later.  Best practices include:
  \begin{itemize}
    \item Use \textbf{self-describing} names (\texttt{release\_date} not
          \texttt{rd} or \texttt{date1}).
    \item Apply a \textbf{consistent naming convention} across the entire
          schema (snake\_case or camelCase -- pick one).
    \item Avoid reserved SQL words as attribute names.
  \end{itemize}
\end{frame}

% ─────────────────────────────────────────────────────────────────────────────
\begin{frame}{Tuples}
  \begin{block}{Definition}
    A \textbf{tuple} is a single row (record) in a relation -- an ordered
    list of values, one per attribute, all conforming to their respective
    domains.  A tuple may contain \texttt{NULL} for optional attributes.
  \end{block}

  \begin{examples}{}
    A tuple of \texttt{Album}:\\[4pt]
    \texttt{(K5MIF8A,\ 4130971072345,\ Fearless,\ Taylor Swift,\ 2008-11-11)}\\[6pt]
    A tuple of \texttt{Track}:\\[4pt]
    \texttt{(T00042,\ K5MIF8A,\ 1,\ 1,\ USUG10801234,\ Love Story,\ 00:03:55)}
  \end{examples}

  The term \emph{tuple} comes from mathematics (a finite ordered list:
  pair, triple, quadruple, \ldots, $n$-tuple).  In relational theory each
  tuple is a complete, self-contained fact about an entity.
  In SQL terminology a tuple is a \textbf{row}; in object-oriented
  terminology it corresponds approximately to an \textbf{instance} of a class.

  \medskip
  A crucial point: within one relation, \textbf{no two tuples may be
  identical} in the strict relational model.  This is why every table
  needs at least one key -- to distinguish tuples from one another.
\end{frame}

% ─────────────────────────────────────────────────────────────────────────────
\begin{frame}{Types of Keys}
  \begin{block}{Overview}
    \textbf{Keys} serve two purposes in the relational model:
    \begin{enumerate}
      \item \textbf{Identification}: uniquely identify every tuple within
            a relation.
      \item \textbf{Reference}: establish relationships between tuples in
            different relations.
    \end{enumerate}
  \end{block}

  \begin{block}{Key Taxonomy}
    \begin{itemize}
      \item \textbf{Super Key} -- any set of attributes that uniquely
            identifies tuples.
      \item \textbf{Candidate Key} -- a minimal super key.
      \item \textbf{Primary Key} -- the chosen candidate key for a relation.
      \item \textbf{Alternate / Secondary Key} -- candidate keys not
            chosen as primary key.
      \item \textbf{Foreign Key} -- an attribute referencing a primary key
            in another relation.
      \item \textbf{Simple Key} -- a key consisting of a single attribute.
      \item \textbf{Composite Key} -- a key consisting of two or more
            attributes.
      \item \textbf{Natural Key} -- a key derived from a real-world attribute.
      \item \textbf{Surrogate Key} -- an artificially generated identifier.
    \end{itemize}
  \end{block}

  Understanding these distinctions is fundamental to good database design and
  is essential for normalisation, referential integrity, and efficient indexing.
\end{frame}

% ─────────────────────────────────────────────────────────────────────────────
\begin{frame}{Super Keys}
  \begin{block}{Definition}
    A \textbf{super key} is \emph{any} set of one or more attributes whose
    combined values are unique across all tuples in a relation.
    \begin{itemize}
      \item A super key may contain \emph{more attributes than necessary}
            for uniqueness -- it can include redundant attributes.
      \item Every relation has at least one super key: the set of
            \emph{all} its attributes is always a super key (trivially).
    \end{itemize}
  \end{block}

  \begin{examples}{}
    Super keys of \texttt{Album(\underline{ID}, EAN, Title, Main\_Artist, Release\_Date)}:\\[4pt]
    $\{$\texttt{ID}$\}$, \quad
    $\{$\texttt{EAN}$\}$, \quad
    $\{$\texttt{ID, EAN}$\}$, \quad
    $\{$\texttt{ID, Title}$\}$, \quad
    $\{$\texttt{ID, EAN, Title}$\}$, \quad \ldots\\[4pt]
    Even $\{$\texttt{ID, EAN, Title, Main\_Artist, Release\_Date}$\}$ is a super key.
  \end{examples}

  The concept of a super key is the \emph{starting point} for the key
  taxonomy.  In practice, super keys with redundant attributes are not
  directly useful -- we refine them into candidate keys by removing
  unnecessary attributes.  However, understanding super keys helps clarify
  \emph{why} minimality matters: smaller keys mean smaller indices, faster
  lookups, and simpler foreign key references.
\end{frame}

% ─────────────────────────────────────────────────────────────────────────────
\begin{frame}{Candidate Keys}
  \begin{block}{Definition}
    A \textbf{candidate key} is a \emph{minimal} super key: no proper
    subset of its attributes is also a super key.
    \begin{itemize}
      \item \textbf{Unique}: no two tuples share the same candidate key value.
      \item \textbf{Irreducible}: removing any attribute from the key breaks
            uniqueness.
      \item A relation can have \textbf{multiple} candidate keys.
    \end{itemize}
  \end{block}

  \begin{examples}{}
    Candidate keys of \texttt{Album}:
    \begin{itemize}
      \item \texttt{ID} -- internal surrogate identifier (UUID or integer)
      \item \texttt{EAN} -- 13-digit product barcode (unique per album)
    \end{itemize}
    Candidate keys of \texttt{Track}:
    \begin{itemize}
      \item \texttt{ID} -- surrogate key
      \item \texttt{ISRC} -- International Standard Recording Code (unique per recording)
      \item $\{$\texttt{Album\_ID, Disc\_Nr, Track\_Nr}$\}$ -- composite natural key
    \end{itemize}
  \end{examples}

  Identifying all candidate keys is a critical step in database design because
  we must choose one as the primary key and enforce the others as unique
  constraints.  Missing a candidate key means potentially allowing
  \emph{logically duplicate} data even when the primary key differs.
\end{frame}

% ─────────────────────────────────────────────────────────────────────────────
\begin{frame}{Primary Key}
  \begin{block}{Definition}
    The \textbf{primary key (PK)} is the candidate key officially chosen
    as the main identifier for tuples in a relation.
    \begin{itemize}
      \item Must be \textbf{unique} for every tuple.
      \item Must \textbf{never be \texttt{NULL}}.
      \item Should be \textbf{stable} -- it should rarely or never change
            after a tuple is created.
    \end{itemize}
  \end{block}

  \begin{examples}{}
    \texttt{Album(\textbf{\underline{ID}}, EAN, Title, Main\_Artist, Release\_Date)}\\
    The primary key is \texttt{ID}; it is underlined by convention.
  \end{examples}

  In SQL, the \texttt{PRIMARY KEY} constraint automatically enforces both
  uniqueness and \texttt{NOT NULL}.  Most database systems also create a
  \textbf{clustered index} on the primary key, which physically orders
  the rows on disk by PK value -- making primary-key lookups extremely fast.

  \begin{alertblock}{Why stability matters}
    If a primary key changes, every foreign key that references it must
    also be updated (or the database must use \texttt{ON UPDATE CASCADE}).
    A changing PK causes cascading updates across potentially millions of
    rows.  This is the main reason surrogate keys are preferred as PKs.
  \end{alertblock}
\end{frame}

% ─────────────────────────────────────────────────────────────────────────────
\begin{frame}{Alternate / Secondary Keys}
  \begin{block}{Definition}
    \textbf{Alternate keys} (also called \emph{secondary keys}) are
    candidate keys that were \emph{not} selected as the primary key.
    They are still genuine unique identifiers -- we simply chose a
    different candidate key to be the PK.
  \end{block}

  \begin{examples}{}
    In \texttt{Album}: Primary key = \texttt{ID},
    so \texttt{EAN} becomes the \textbf{alternate key}.\\[4pt]
    In SQL: \quad \texttt{UNIQUE (EAN)} \quad -- enforces uniqueness without
    making \texttt{EAN} the primary key.
  \end{examples}

  Alternate keys serve an important practical role: \textbf{external systems}
  often know a natural identifier (EAN, ISBN, ISRC) but not the internal
  surrogate ID.  By enforcing a unique constraint on the alternate key, you
  allow efficient lookups from those external systems while keeping a
  clean surrogate PK for internal foreign key references.

  \medskip
  A good database schema \textbf{enforces all candidate keys}, not just
  the primary key.  Failing to do so allows logically duplicate records
  that differ only in their surrogate ID -- a common source of data
  quality problems in production systems.
\end{frame}

% ─────────────────────────────────────────────────────────────────────────────
\begin{frame}{Foreign Keys}
  \begin{block}{Definition}
    A \textbf{foreign key (FK)} is an attribute (or set of attributes)
    in one relation whose values must match an existing primary key value
    in another (or the same) relation -- or be \texttt{NULL}.
    \begin{itemize}
      \item Enforces \textbf{referential integrity}: you cannot reference
            a tuple that does not exist.
      \item Creates a \textbf{parent--child} relationship between tables.
    \end{itemize}
  \end{block}

  \begin{examples}{}
    \texttt{Track(\underline{ID}, \textit{Album\_ID}\textsuperscript{FK}, Disc\_Nr, \ldots)}\\[4pt]
    \texttt{Album\_ID} in \texttt{Track} references \texttt{ID} in \texttt{Album}.\\
    Inserting a track for a non-existent album raises an integrity error.
  \end{examples}

  \begin{alertblock}{Referential Integrity Actions}
    When a referenced tuple is deleted or its PK is updated, the database
    must decide what to do with referencing rows.  SQL offers:
    \texttt{RESTRICT}, \texttt{CASCADE}, \texttt{SET NULL}, \texttt{SET DEFAULT}.
    Choosing the right action is part of schema design -- \texttt{CASCADE DELETE}
    can silently destroy large amounts of data if misapplied.
  \end{alertblock}
\end{frame}

% ─────────────────────────────────────────────────────────────────────────────
\begin{frame}{Simple vs.\ Composite Keys}
  \begin{block}{Simple Key}
    A \textbf{simple key} consists of a \textbf{single attribute}.
    \begin{itemize}
      \item Easy to reference as a foreign key (one column).
      \item Typically a surrogate (auto-increment integer or UUID).
    \end{itemize}
  \end{block}

  \begin{block}{Composite Key}
    A \textbf{composite key} consists of \textbf{two or more attributes}
    whose combined values are unique.
    \begin{itemize}
      \item Common in \textbf{junction / association tables} that
            represent many-to-many relationships.
      \item Natural uniqueness constraints often emerge as composite keys.
    \end{itemize}
  \end{block}

  \begin{examples}{}
    Simple key: \texttt{ID} in \texttt{Album}\\[4pt]
    Composite key: \texttt{\{Album\_ID, Disc\_Nr, Track\_Nr\}} in \texttt{Track}\\
    $\Rightarrow$ Disc 1, Track 3 of album ``K5MIF8A'' identifies exactly one track.
  \end{examples}

  \textbf{Rule of thumb}: prefer a simple surrogate primary key.
  Composite keys as PKs make foreign key references more verbose (all
  columns must be replicated) and join conditions more complex.
  Use composite keys as \emph{unique constraints} to enforce natural
  business rules, but keep the PK simple.
\end{frame}

% ─────────────────────────────────────────────────────────────────────────────
\begin{frame}{Natural vs.\ Surrogate Keys}
  \begin{block}{Natural Key}
    A \textbf{natural key} is derived from a real-world attribute that
    already exists in the data domain.
    \begin{itemize}
      \item Examples: EAN (barcode), ISBN (book), ISRC (music track),
            social security number, email address.
      \item \textbf{Pros}: human-readable, meaningful, no artificial column needed.
      \item \textbf{Cons}: may change (email addresses change!), often long
            (bad for indices and FK columns), may be unavailable at insert time.
    \end{itemize}
  \end{block}

  \begin{block}{Surrogate Key}
    A \textbf{surrogate key} is an artificially generated identifier with
    no inherent business meaning.
    \begin{itemize}
      \item Examples: auto-increment integer (\texttt{SERIAL}), UUID v4.
      \item \textbf{Pros}: stable, compact, always available, never changes.
      \item \textbf{Cons}: meaningless on its own, requires joining to see
            human-readable identity.
    \end{itemize}
  \end{block}

  \begin{alertblock}{Best Practice}
    Use a \textbf{surrogate key as the primary key} and enforce natural
    keys as \texttt{UNIQUE} constraints (alternate keys).  You get the
    stability of a surrogate and the integrity guarantee of the natural key.
  \end{alertblock}
\end{frame}

% ─────────────────────────────────────────────────────────────────────────────
\begin{frame}{Exercise: Relational Databases}
  \begin{block}{Task}
    Work through relations \textbf{Exercise 1A} through \textbf{Exercise 1D}.\\
    (Simple: 1A \& 1B \quad Intermediate: 1C \& 1D)
  \end{block}

  For each relation, answer the following questions:

  \begin{enumerate}
    \item List all \textbf{attributes} and their domains.
    \item Identify all \textbf{candidate keys}; select and justify your choice of
          \textbf{primary key}.
    \item Identify all \textbf{foreign keys} and state which primary key each
          one references.
    \item Classify every key as
          \emph{simple / composite} and \emph{natural / surrogate}.
    \item Create \textbf{2 realistic sample tuples} for each relation that
          respect all constraints.
  \end{enumerate}

  \begin{alertblock}{Time: 15 minutes}
    Then we will discuss the solutions together.  There is often more than one
    valid choice of primary key -- be prepared to justify yours.
  \end{alertblock}
\end{frame}

% ─────────────────────────────────────────────────────────────────────────────
\begin{frame}{Preview: Relational Algebra}
  \begin{block}{Why Do We Need It?}
    Defining relations, attributes, and keys tells us \emph{how to store}
    data.  We now need a formal language to describe \emph{how to query
    and manipulate} that data.
  \end{block}

  \textbf{Relational algebra} provides the mathematical foundation for all
  relational query languages, including SQL.  It defines a small set of
  \textbf{operations on relations that produce new relations} (closure property),
  which can be freely composed.

  \medskip
  Core operations we will cover:

  \begin{itemize}
    \item Selection $\sigma$ -- filter rows by a predicate
    \item Projection $\pi$ -- select a subset of columns
    \item Rename $\rho$ -- rename a relation or attribute
    \item Union $\cup$, Intersection $\cap$, Difference $\setminus$ -- set operations
    \item Cartesian product $\times$ -- combine every pair of tuples
    \item Join $\bowtie$ -- combine related tuples (natural join)
    \item Left / Right outer join $\ltimes$ / $\rtimes$
    \item Division $\div$ -- ``find all X that are associated with all Y''
  \end{itemize}

  Every SQL query you will ever write is an expression in relational algebra
  in disguise -- understanding the algebra makes SQL comprehensible rather
  than a collection of memorised syntax.
\end{frame}

% ─────────────────────────────────────────────────────────────────────────────
\begin{frame}{Preview: Normal Forms}
  \begin{block}{Motivation}
    Not all relation designs are equally good.  A poorly designed schema
    can lead to:
    \begin{itemize}
      \item \textbf{Redundancy} -- the same fact stored in multiple places.
      \item \textbf{Update anomalies} -- changing one fact requires updating
            many rows; missing one causes inconsistency.
      \item \textbf{Insertion anomalies} -- you cannot record a fact without
            first recording an unrelated fact.
      \item \textbf{Deletion anomalies} -- deleting one fact accidentally
            destroys another.
    \end{itemize}
  \end{block}

  \begin{block}{Normal Forms}
    \textbf{Normal forms} are a hierarchy of design rules that systematically
    eliminate these problems by ensuring each fact is stored exactly once.
    We will cover:
    \begin{itemize}
      \item \textbf{1NF} -- atomic values, no repeating groups
      \item \textbf{2NF} -- no partial functional dependencies on a composite PK
      \item \textbf{3NF} -- no transitive functional dependencies
      \item \textbf{BCNF} -- Boyce--Codd Normal Form (stronger 3NF)
      \item \textbf{4NF / 5NF} -- optional, for multi-valued and join dependencies
    \end{itemize}
  \end{block}

  Normalisation is the \emph{bridge} between a rough data model and a
  production-quality database schema.
\end{frame}
